\section{Related Work and Future Directions}
\label{sec:relatedWork}


In this paper, we presented the first logical relation for session types
which supports GV-style rules, can prove deadlock freedom, and
admits safe processes which do not comply with communication topology
constraints imposed by the syntactic type system.
In this section we review the literature in view of the technical development
and sketch potential future developments.

\citet{DBLP:journals/iandc/PerezCPT14} were the first to investigate
session types using logical relations.
Their starting point is the logical account of session types \cite{DBLP:conf/concur/CairesP10}; their process calculus includes a (non-recursive)
treatment of the exponential, and uses the logical relation to prove
strong normalization, confluence, and behavioral equality.
More recent work by \citet{conf/esop/RochaC23} uses a logical relation that specifies
strong normalization and safe handling of state for message-passing
processes using linear session types
(based on an extension of ILL~\cite{conf/tapsoft/GirardL87}).
\mw{claude says this should be: (based on an extension of second-order linear logic~\cite{DBLP:journals/tcs/Girard87})}
In both \cite{DBLP:journals/iandc/PerezCPT14,conf/esop/RochaC23},
the logical relation is defined only for syntactically typed processes.
Deadlock freedom is established by proving subject reduction and progress
on the syntactic derivations, and the proofs of the logical relation
rely on having externally established deadlock freedom.
In contrast, we establish all the behavioral guarantees of the type system
as a consequence of the semantic logical relation.

An interesting direction for further work would be to establish
compositional and liberal logical relations for properties beyond
deadlock freedom.
For example,
\citet{conf/esop/Atkey17} defines an operational semantics for
Wadler's ``processes as CLL proofs''~\cite{conf/icfp/Wadler12} and
uses a logical relation to prove adequacy
of the operational semantics and a type-based
denotational semantics with the goal
of establishing observational equivalences.
\citet{DBLP:conf/lics/DerakhshanBJ21} study noninterference for
information flow control.
Van den Heuvel et al.~\cite{DBLP:conf/ecoop/0001DB24} extend the
work to cyclic processes.
\citet{DBLP:journals/pacmpl/YaoICBZ25} consider timed
message passing for IoT systems via session types defined by a logical relation.

The logical-relations approach also offers a nice opportunity
when extending the language of types to support polymorphism,
as it would allow the derivation of parametricity theorems from polymorphic types.


Another direction would be to extend the features of the programming model.
For example, LinearActris~\cite{JacobsHK24} supports mutable heaps
in addition to message passing primitives.
It would be interesting to try and embed our logical relation into
a separation logic, allowing a more flexible mix of paradigms.

While this work focused on the Curry-Howard based ways of composing processes,
approaches based on priorities, \eg \cite{Kobayashi97,Padovani14,DardhaG18,KokkeD23,APCP},
track dependencies between actions across sessions explicitly in the types.
Finding a semantic logical relation capturing priority-based reasoning
would provide powerful compositional means of proving well-behavedness of
components that cannot be syntactically typed (like our load balancer example).
It could also provide a semantic bridge between priority-based and Curry-Howard-based systems, in the spirit of \cite{DardhaP22,Jaramillo025}.

Multiparty Session Types~\cite{HondaYC16} also shed light on global dependencies in types.
For programs interleaving several multiparty sessions,
establishing deadlock freedom requires additional priority-based annotations and checks.
If one considers programs with a single session involving multiple participants,
the Multiparty Session Type discipline ensures deadlock freedom.
The semantic logical relation understanding of Multiparty Session Types is
still underdeveloped.
The Multris logic~\cite{HinrichsenJK24} made some steps in developing
a logical relation embedded into the Iris separation logic,
but it can only prove safe usage of channels and resources,
not deadlock freedom.

Finally, a challenging open problem is to provide a semantic
and compositional logical relation to prove \emph{liveness} properties
like lock freedom.
Priority-based systems have been adapted to prove lock freedom and progress
properties~\cite{Kobayashi02,Padovani13},
but the soundness argument is syntactic.
Several works in Concurrent Separation Logics 
\cite{DOsualdoSFG21,LiangF18,LeeLYCKH25}
explored compositional semantic characterizations of progress properties
for shared memory concurrent programs.
Combining those ideas with our logical relation approach
could provide some leverage to attack the problem.

%%% Local Variables:
%%% mode: latex
%%% TeX-master: "main"
%%% End:
